\chapter{问题集}
\label{chap:question}
提出问题比解决问题更重要，因为提出问题需要更深入的理解和创造力。以下是一些关于量子场论、规范场论以及一些奇妙想法的问题，不需要回答，可以留着慢慢思考，并最终体现于博士学位论文中。
\section{量子场论}
\begin{itemize}
    \item 在QED中，通过讨论电子自能修正过程，引入了重整化的概念。比如on-shell重整化，我们要求修正后的电子质量等于物理质量，修正后的电荷等于物理电荷，传播子在pole处的留数等于1，而且在微扰论的逐阶都是well-defined。那为什么在类似的重夸克自能修正过程中，会在On-shell重整化方案下出现renormalon的现象呢，这背后反映了 on-shell重整化方案的什么样的的局限性。renormalon的出现与否是否具有重整化方案的依赖？
    \item 质量作为QCD理论中的一个naive的参数，从有效理论的角度出发，应该被作为一个Wilosn系数来处理。这样看应该具备scheme dependence，以及关于能标的跑动。在不同的有效理论中会吸收来自短程的微扰贡献，这一观点是如何体现的，比如pNRQCD（potential NRQCD），HQET等等。
\[
\Sigma_{\rm QCD}(\slashed p)\sim g_s^2 C_F
\int \frac{d^4k}{(2\pi)^4}\,
\gamma^\mu \frac{i(\slashed p-\slashed k+m)}{(p-k)^2-m^2+i0}\gamma^\nu
\frac{-ig_{\mu\nu}}{k^2+i0},
\]
\end{itemize}
\section{规范场论}
\section{奇妙想法}